\subsection{Independence and Compatibility of the Axioms}
\label{sec:ch3-2-independence}

There is one class of field theories known to satisfy axioms O, I, II,
III, and IV: the theories of free fields of various masses and spins,
described heuristically at the end of Section \ref{sec:1.4}.  These theories will
now be described more precisely.  That they exist shows that the axioms are
compatible, but in a physically trivial way since the particles they
describe do not interact.

In all the free field theories the total number of particles is an integral
of motion and the Hilbert space of states is written as a direct sum,
\begin{equation}
  \Hilbert=\bigoplus_{n=0}^{\infty}\Hilbert^{(n)},
  \tag{3-9}\label{eq:ch3-free-field-hilbert-sum}
\end{equation}
where $\Hilbert^{(n)}$ is the subspace of those states with exactly $n$
particles.  This means that a vector, $\ket{\Phi}$, of $\Hilbert$ is given
by a sequence
\begin{equation}
  \{\Phi^{(0)},\Phi^{(1)},\Phi^{(2)},\ldots\}
  \tag{3-10}\label{eq:ch3-free-field-sequence}
\end{equation}
of vectors with $\Phi^{(j)}\in\Hilbert^{(j)}$, and the scalar product in
$\Hilbert$ is
\begin{equation}
  \braket{\Phi}{\Psi}=\sum_{n=0}^{\infty}
  \braket{\Phi^{(n)}}{\Psi^{(n)}},
  \tag{3-11}\label{eq:ch3-free-field-inner-product}
\end{equation}
where the scalar product $\braket{\Phi^{(n)}}{\Psi^{(n)}}$ is in
$\Hilbert^{(n)}$.  Only those sequences are in $\Hilbert$ which satisfy
\begin{equation}
  \braket{\Phi}{\Phi}=\sum_{n=0}^{\infty}\norm{\Phi^{(n)}}^2<\infty.
  \tag{3-12}\label{eq:ch3-free-field-norm}
\end{equation}

For all spins, $\Hilbert^{(0)}$ is the one-dimensional Hilbert space of
complex numbers.  For spin $s$, the Hilbert space $\Hilbert^{(n)}$ is the
space of all square-integrable functions of the arguments
$p_1\alpha_1,\ldots,p_n\alpha_n$, symmetric if $s$ is integral and
anti-symmetric if $s$ is half-odd-integral.  Here $\alpha_j$ stands for a
group of undotted spinor indices like those in \eqref{eq:1-58}.  The scalar product
in $\Hilbert^{(n)}$ is
\begin{equation}
\begin{aligned}
  \braket{\Phi^{(n)}}{\Psi^{(n)}}={}&
  \int\!\cdots\!\int
  \dd\Omega_m(p_1)\cdots\dd\Omega_m(p_n)
  \sum_{\substack{\alpha_1\ldots\alpha_n\\\beta_1\ldots\beta_n}}
  \Phi^{(n)*}(p_1\alpha_1,\ldots,p_n\alpha_n) \\
  &\quad\times\prod_{j=1}^{n}
  \mathcal{D}^{(s,0)}\!\left(\frac{\widetilde p_j}{m}\right)_{\alpha_j\beta_j}
  \Psi^{(n)}(p_1\beta_1,\ldots,p_n\beta_n).
\end{aligned}
\tag{3-13}\label{eq:ch3-free-field-n-particle-product}
\end{equation}
and the field operator $\varphi_\alpha$ is defined by
\begin{equation}
\begin{aligned}
  \varphi_\alpha(f)\Psi^{(n)}(p_1\alpha_1,\ldots,p_n\alpha_n)
  ={ }&\sqrt{\pi}\Biggl\{
  \sqrt{n+1}\int\!\dd\Omega_m(p)\,f(p)
  \Psi^{(n+1)}(p\alpha,p_1\alpha_1,\ldots,p_n\alpha_n) \\
  &\quad+\frac{1}{\sqrt n}\sum_{j=1}^{n}(-1)^{2s(j+1)}
  \widetilde f(-p_j)\mathcal{D}^{(s,0)}(\zeta)_{\alpha\alpha_j} \\
  &\qquad\times\Psi^{(n-1)}(p_1\alpha_1,\ldots,
  \widehat p_j\widehat\alpha_j,\ldots,p_n\alpha_n)\Biggr\}.
\end{aligned}
\tag{3-14}\label{eq:ch3-free-field-operator}
\end{equation}

Here a hat over a letter means omit it, and $\widetilde f$ is defined by
\begin{equation}
  \widetilde f(p)=\frac{1}{(2\pi)^2}\int
  \ee^{+\ii p\cdot x}f(x)\,\dd x.
  \tag{3-15}\label{eq:ch3-fourier-transform}
\end{equation}

The representation of the inhomogeneous $SL(2,\mathbb C)$ is here
\begin{equation}
\begin{aligned}
  \bigl(U(a,A)\Psi\bigr)^{(n)}(p_1\alpha_1,\ldots,p_n\alpha_n)
  ={ }&\exp\!\left[-\ii\left(\sum_{j=1}^{n}p_j\right)\!\cdot a\right]
  \sum_{(\beta)}
  \prod_{j=1}^{n}\mathcal{D}^{(s,0)}(A)_{\alpha_j\beta_j} \\
  &\quad\times\Psi^{(n)}\bigl(\Lambda(A^{-1})p_1\beta_1,\ldots,
  \Lambda(A^{-1})p_n\beta_n\bigr).
\end{aligned}
\tag{3-16}\label{eq:ch3-inhomogeneous-representation}
\end{equation}

For the domain $D$, one can take, for example, those $\ket{\Psi}$ for
which $\Psi^{(n)}$ vanishes for all sufficiently large $n$ and is the
restriction to the hyperboloids $p_j^2=-m^2$ of an infinitely differentiable
function in $\Schwartz(\mathbb R^{4n})$.  It is a straightforward task to
verify O, I, II, III, and IV for the theory given by these definitions, an
exercise we strongly recommend to the reader.  For the case of spin zero,
the representation of the inhomogeneous $SL(2,\mathbb C)$ is just that
illustrated in Figure \ref{fig:1-3}.

There are fields different from free fields, which can be constructed in
terms of free fields.  For example, consider the fields which are indicated
in terms of the scalar field, $\varphi$, by the limiting procedure
\begin{equation}
\begin{aligned}
  :\!D^\alpha\varphi(x)D^\beta\varphi(x)\!:
  ={ }&\lim_{x_1,x_2\to x}\Bigl[
  D^\alpha\varphi(x_1)D^\beta\varphi(x_2)
  -\bra{\Omega}D^\alpha\varphi(x_1)D^\beta\varphi(x_2)\ket{\Omega}\Bigr] \\
  &\text{and} \\
  :\!D^\alpha\varphi(x)D^\beta\varphi(x)D^\gamma\varphi(x)\!:
  ={ }&\lim_{x_1,x_2,x_3\to x}\Bigl[
  D^\alpha\varphi(x_1)D^\beta\varphi(x_2)D^\gamma\varphi(x_3) \\
  &\quad-\bra{\Omega}D^\alpha\varphi(x_1)D^\beta\varphi(x_2)\ket{\Omega}
  D^\gamma\varphi(x_3) \\
  &\quad-\bra{\Omega}D^\alpha\varphi(x_1)D^\gamma\varphi(x_3)\ket{\Omega}
  D^\beta\varphi(x_2) \\
  &\quad-\bra{\Omega}D^\beta\varphi(x_2)D^\gamma\varphi(x_3)\ket{\Omega}
  D^\alpha\varphi(x_1)\Bigr].
\end{aligned}
\tag{3-17}\label{eq:ch3-wick-polynomials}
\end{equation}
and analogously for higher degrees in $\varphi$.  [Recall the definition of
$D^\alpha$ in \eqref{eq:ch2-derivative-multiindex}.]  It is not difficult to show, by using formula
\eqref{eq:ch3-free-field-operator}, that the right-hand sides of these expressions indeed define fields
satisfying O, I, II, and III with the same domain $D$ as that of the free
field, and the same representation $U(a,A)$.  They have transformation laws
appropriate to the tensor indices $\alpha,\beta$, or $\alpha,\beta,\gamma$,
etc., respectively.  By contracting the indices we get scalar fields, for
example,
\begin{equation}
  \psi(x)=\varphi(x)+:
  \frac{\partial}{\partial x_\mu}\varphi(x)
  \frac{\partial}{\partial x^\mu}\varphi(x): .
  \tag{3-18}\label{eq:ch3-wick-scalar-example}
\end{equation}

Such fields are sometimes called Wick polynomials.  The particular example
\eqref{eq:ch3-wick-scalar-example} satisfies axiom IV in addition to O, I, II, and III, but it gives the
same $S$-matrix as the free field.  That is what always happens for a field
theory defined by a Wick polynomial, as will be seen in Chapter \ref{ch:general-theorems-relativistic-qft}.

Another class of examples of fields can be obtained by replacing the
$\dd\Omega_m(p)$ in \eqref{eq:ch3-free-field-n-particle-product} and \eqref{eq:ch3-free-field-operator} by
$\int_0^\infty\rho(a)\,\dd a\,\dd\Omega_a(p)$, where $\rho(a)$ is a
non-negative weight.  Unless $\rho(a)=\delta(a-m)$, this gives a
representation of the Lorentz group different from that associated with the
free field of mass $m$.  Such fields are called \emph{generalized free
fields}.  It is proved in Ref. 4 that axiom IV is not satisfied for certain
choices of $\rho$.  This shows that IV is independent of O, I, II, and III.

These examples show that there are field theories whose fields satisfy O, I,
II, and III for certain representations $\{a,A\}\longrightarrow U(a,A)$.
Examples of representations also exist which have no non-trivial fields
having the vacuum as cyclic vector.  For example, that is the case if the
representation contains no masses above some limit.  This is a consequence
of a theorem proved in Ref. 15, which says that the energy-momentum spectrum
of a field theory must be additive; that is, if $p_1$ and $p_2$ are in the
spectrum so is $p_1+p_2$.

It is fairly obvious that III is independent of O, I, and II since it is easy
to write down non-local fields satisfying O, I, and II.  In fact it has been
repeatedly suggested that local commutativity (axiom III) is too strong an
assumption, since there is no evidence for (or against) the simultaneous
measurability of fields at points of space-time which are separated by very
small distances, say $10^{-16}\,\mathrm{cm}$.  In this connection Theorem \ref{thm:ch4-global-locality}
is relevant.  If III is replaced by commutativity at large space-like
distances, say $(x-y)^2>\ell^2$ (thus introducing a fundamental length), then
III follows.  Thus if III is to be essentially weakened the commutator must
be allowed to be different from zero at all space-like distances.

It has been suggested that the irreducibility assumption in the definition
of a field theory should be strengthened by the \emph{time slice axiom} as
follows: if $T$ is a set in space-time lying between two space-like surfaces,
then operators of the form $\varphi(f)$ should form an irreducible set,
where $f=0$ outside $T$.  This is a mathematically precise substitute for
the ill-defined requirement of irreducibility of the fields and their
derivatives at a given time.  We shall not analyze the time slice axiom here.
We note that it is stronger than assuming the cyclicity of the vacuum, as
has been shown in Ref. 3.

The examples given in this section are trivial in the sense that, though they
satisfy O, I, II, and III, they yield a trivial collision theory.
